% GENERATED FILE. Edit canonical lessons, then run npm run generate:books.
% canonical-source-hash: fnv1a64:05cd652c9e8811ad
% canonical-lessons: LA-C49-possum, LA-C49-volo, LA-C49-qui, LA-C49-et, LA-C49-quando

\chapter{Possum, Volo, Qui, Et, Quando}
\label{ch:la-formae}
\hlchaptermodality{49}

\begin{chapteropening}
\textbf{By the end of this chapter:} \emph{I can say what I can do and what I want, join two statements, and put a when-clause or a who-clause after a noun.}

The fifth word of the group, taken apart against the four before it.
\end{chapteropening}

\section[possum, posse]{possum — potis plus sum, and the verb that will not sit still}

\label{lesson:LA-C49-possum}



\begin{quote}
Six recalls before the new word, at growing distances. They are not decoration: each one is the retrieval window this book measures, and each names a word far enough back that saying it is work.

\begin{itemize}
\raggedright
  \item \emph{Recall:} say \emph{aut} — \textbf{R1}, one lesson back, before it decays
  \item \emph{Recall:} say \emph{nesciō} — \textbf{R2}, the first real retrieval, five to fifteen lessons back
  \item \emph{Recall:} say \emph{salvē (post merīdiem)} — \textbf{R3}, durable, twenty to sixty lessons back
  \item \emph{Recall:} say \emph{salvē / salvēte} — \textbf{R3}, durable again, a second word from the same distance
  \item \emph{Recall:} say \emph{pānis} — \textbf{R4}, recognition at distance, eighty lessons back or more
  \item \emph{Recall:} say \emph{Iānuārius} — \textbf{R4}, recognition at distance, a second word from that far back
\end{itemize}
\end{quote}

\begin{cousinweb}
\textbf{possum} — \textquotedblleft{}\textbf{I am able, I can}.\textquotedblright{}

Two words in a coat. \textbf{Potis} meant \textquotedblleft{}able,\textquotedblright{} and Latin bolted it onto \textbf{sum}, the verb the eight-verb chapter gave you. Where \emph{sum} begins with \emph{s}, the \emph{t} of \emph{potis} flattens into it — \emph{pot-sum} becomes \textbf{possum} — and where it does not, the \emph{t} survives:

\noindent
\begin{tabularx}{\linewidth}{@{}>{\raggedright\arraybackslash}X>{\raggedright\arraybackslash}X>{\raggedright\arraybackslash}X@{}}
\toprule
\textbf{person} & \textbf{form} & \textbf{built from} \\
\midrule
I can & \textbf{possum} & \emph{pot-} + \emph{sum} \\
you can & \textbf{potes} & \emph{pot-} + \emph{es} \\
he/she can & \textbf{potest} & \emph{pot-} + \emph{est} \\
we can & \textbf{possumus} & \emph{pot-} + \emph{sumus} \\
they can & \textbf{possunt} & \emph{pot-} + \emph{sunt} \\
\bottomrule
\end{tabularx}

What follows \emph{possum} is an \textbf{infinitive}, which is the form every verb in this book is already stored under: \textbf{possum vidēre}, \textquotedblleft{}I can see.\textquotedblright{} \textbf{Nōn possum vidēre}, \textquotedblleft{}I cannot see,\textquotedblright{} using last chapter's rule.

English \textbf{potent}, \textbf{possible}, \textbf{potential} and \textbf{omnipotent} are all \emph{potis}; so is \textbf{power}, by way of French \emph{pouvoir}.
\end{cousinweb}

\subsection*{Your turn}
Say these aloud:

\begin{itemize}
\raggedright
  \item \emph{possum, posse}
  \item every word this chapter has given you so far, in order
\end{itemize}

\begin{tcolorbox}[breakable,colback=teal!4,colframe=teal!35!black,title={Before you move on}]
What two words are inside \emph{possum}? (\emph{\textbf{Potis}} and \emph{\textbf{sum}}.) What form follows it? (\textbf{The infinitive}.) Say \textquotedblleft{}I cannot see.\textquotedblright{} (\emph{\textbf{Nōn possum vidēre}}.)
\end{tcolorbox}
\section[volō, velle]{volō — and the root inside voluntary, benevolent and will itself}

\label{lesson:LA-C49-volo}



\begin{quote}
Six recalls before the new word, at growing distances. They are not decoration: each one is the retrieval window this book measures, and each names a word far enough back that saying it is work.

\begin{itemize}
\raggedright
  \item \emph{Recall:} say \emph{possum} — \textbf{R1}, one lesson back, before it decays
  \item \emph{Recall:} say \emph{male} — \textbf{R2}, the first real retrieval, five to fifteen lessons back
  \item \emph{Recall:} say \emph{salvē — bonum vesperum — bonam noctem — valē} — \textbf{R3}, durable, twenty to sixty lessons back
  \item \emph{Recall:} say \emph{sum} — \textbf{R3}, durable again, a second word from the same distance
  \item \emph{Recall:} say \emph{medius diēs} — \textbf{R4}, recognition at distance, eighty lessons back or more
  \item \emph{Recall:} say \emph{hōra} — \textbf{R4}, recognition at distance, a second word from that far back
\end{itemize}
\end{quote}

\begin{cousinweb}
\textbf{volō} — \textquotedblleft{}\textbf{I want, I wish}.\textquotedblright{}

The word that turns everything you own into a request. Like \emph{possum}, it takes an \textbf{infinitive}:

\begin{quote}
\textbf{Volō bibere.} — I want to drink. \textbf{Nōn volō dormīre.} — I do not want to sleep.
\end{quote}

Its forms are irregular and worth hearing once: \textbf{volō, vīs, vult, volumus, vultis, volunt}. The infinitive is \textbf{velle}, not \emph{volere} — an old form that survived because the verb was too common to regularise.

The root is Proto-Indo-European \emph{\textbf{*welh\textsubscript{1}-}}, \textquotedblleft{}to choose.\textquotedblright{} English inherited it directly as \textbf{will} — a blood cousin, not a loan — and borrowed it back from Latin many times over: \textbf{voluntary}, \textbf{volition}, \textbf{benevolent} (\textquotedblleft{}well-wishing\textquotedblright{}), \textbf{malevolent} (with last chapter's \emph{male}), and \textbf{voluptuous}.

And you have already met its imperative. \textbf{Vel}, \textquotedblleft{}or,\textquotedblright{} from the last lesson of the denying chapter, is literally \emph{volō} telling you to choose.
\end{cousinweb}

\subsection*{Your turn}
Say these aloud:

\begin{itemize}
\raggedright
  \item \emph{volō, velle}
  \item every word this chapter has given you so far, in order
\end{itemize}

\begin{tcolorbox}[breakable,colback=teal!4,colframe=teal!35!black,title={Before you move on}]
What form follows \emph{volō}? (\textbf{The infinitive}.) Which English word is a blood cousin rather than a loan? (\emph{\textbf{Will}}.) And where have you already met \emph{volō}'s imperative? (\textbf{In \emph{vel}}, \textquotedblleft{}or.\textquotedblright{})
\end{tcolorbox}
\section[quī, quae, quod]{quī — the commonest word in Latin prose, and the one that joins two sentences into one}

\label{lesson:LA-C49-qui}



\begin{quote}
Six recalls before the new word, at growing distances. They are not decoration: each one is the retrieval window this book measures, and each names a word far enough back that saying it is work.

\begin{itemize}
\raggedright
  \item \emph{Recall:} say \emph{volō} — \textbf{R1}, one lesson back, before it decays
  \item \emph{Recall:} say \emph{malus} — \textbf{R2}, the first real retrieval, five to fifteen lessons back
  \item \emph{Recall:} say \emph{sum} — \textbf{R3}, durable, twenty to sixty lessons back
  \item \emph{Recall:} say \emph{habeō} — \textbf{R3}, durable again, a second word from the same distance
  \item \emph{Recall:} say \emph{medius diēs} — \textbf{R4}, recognition at distance, eighty lessons back or more
  \item \emph{Recall:} say \emph{hōra} — \textbf{R4}, recognition at distance, a second word from that far back
\end{itemize}
\end{quote}

\begin{cousinweb}
\textbf{quī, quae, quod} — \textquotedblleft{}\textbf{who, which, that}.\textquotedblright{}

This is the single most common word in surviving Latin prose, and it does one job: it hangs a second statement onto a noun in the first.

\begin{quote}
\textbf{Amīcus quī venit.} — The friend \textbf{who} comes. \textbf{Liber quem legō.} — The book \textbf{which} I am reading.
\end{quote}

Three forms, because it agrees with the noun it hangs on: \textbf{quī} for a masculine one, \textbf{quae} for a feminine one, \textbf{quod} for a neuter one.

You already know its relatives. \textbf{Quid} (\textquotedblleft{}what\textquotedblright{}) from the wellbeing exchange and \textbf{quis} (\textquotedblleft{}who\textquotedblright{}) are the question words built on the same Proto-Indo-European stem \emph{\textbf{*kwi-}}. And Grimm's Law — the \emph{piscis} law, \emph{p} to \emph{f} — does \emph{kw} to \emph{hw} as well, which is why English \textbf{who}, \textbf{what}, \textbf{which} and \textbf{why} are inherited cousins of this word rather than borrowings from it. The spelling \emph{wh-} is the old \emph{hw-} with its letters swapped.
\end{cousinweb}

\subsection*{Your turn}
Say these aloud:

\begin{itemize}
\raggedright
  \item \emph{quī, quae, quod}
  \item every word this chapter has given you so far, in order
\end{itemize}

\begin{tcolorbox}[breakable,colback=teal!4,colframe=teal!35!black,title={Before you move on}]
What does \emph{quī} do to two sentences? (\textbf{Joins them} — hangs the second on a noun in the first.) Are English \emph{who} and \emph{what} cousins or loans? (\textbf{Cousins}, by Grimm's Law.)
\end{tcolorbox}
\section[et]{et — and the enclitic -que, which does the same job glued to the back of a word}

\label{lesson:LA-C49-et}



\begin{quote}
Six recalls before the new word, at growing distances. They are not decoration: each one is the retrieval window this book measures, and each names a word far enough back that saying it is work.

\begin{itemize}
\raggedright
  \item \emph{Recall:} say \emph{quī} — \textbf{R1}, one lesson back, before it decays
  \item \emph{Recall:} say \emph{sed} — \textbf{R2}, the first real retrieval, five to fifteen lessons back
  \item \emph{Recall:} say \emph{caelum} — \textbf{R3}, durable, twenty to sixty lessons back
  \item \emph{Recall:} say \emph{eō} — \textbf{R3}, durable again, a second word from the same distance
  \item \emph{Recall:} say \emph{hōra} — \textbf{R4}, recognition at distance, eighty lessons back or more
  \item \emph{Recall:} say \emph{vīgintī annōs nātus sum} — \textbf{R4}, recognition at distance, a second word from that far back
\end{itemize}
\end{quote}

\begin{cousinweb}
\textbf{et} — \textquotedblleft{}\textbf{and}.\textquotedblright{}

The plainest word in this book, and the corpus went forty-seven chapters without teaching it.

\begin{quote}
\textbf{Aqua et vīnum.} — Water and wine. \textbf{Pater et māter.} — Father and mother.
\end{quote}

Latin has a \textbf{second} way of saying it, and a reader meets it constantly: \textbf{-que}, glued onto the \textbf{back} of the second word rather than standing between the two.

\begin{quote}
\textbf{Senātus populusque Rōmānus} — \textquotedblleft{}The Senate \textbf{and} the Roman people,\textquotedblright{} the SPQR on every Roman standard. The \emph{-que} is that fourth letter.
\end{quote}

\emph{-que} is Proto-Indo-European \emph{\textbf{*kwe}}, an enclitic Sanskrit kept as \emph{ca} and Greek as \emph{te}. \textbf{Et} went the other way: it is Spanish \emph{y}, French \emph{et}, Italian \emph{e} and Portuguese \emph{e}, and every Romance daughter kept it.
\end{cousinweb}

\subsection*{Your turn}
Say these aloud:

\begin{itemize}
\raggedright
  \item \emph{et}
  \item every word this chapter has given you so far, in order
\end{itemize}

\begin{tcolorbox}[breakable,colback=teal!4,colframe=teal!35!black,title={Before you move on}]
Two ways to say \emph{and} in Latin? (\emph{\textbf{Et}}, and \textbf{-que} glued to the back.) What does the \emph{Q} in SPQR stand for? (\emph{\textbf{-que}}, \textquotedblleft{}and.\textquotedblright{})
\end{tcolorbox}
\section[quandō]{quandō — as a question and as the hinge of a clause}

\label{lesson:LA-C49-quando}



\begin{quote}
Six recalls before the new word, at growing distances. They are not decoration: each one is the retrieval window this book measures, and each names a word far enough back that saying it is work.

\begin{itemize}
\raggedright
  \item \emph{Recall:} say \emph{et} — \textbf{R1}, one lesson back, before it decays
  \item \emph{Recall:} say \emph{aut} — \textbf{R2}, the first real retrieval, five to fifteen lessons back
  \item \emph{Recall:} say \emph{mare} — \textbf{R3}, durable, twenty to sixty lessons back
  \item \emph{Recall:} say \emph{veniō} — \textbf{R3}, durable again, a second word from the same distance
  \item \emph{Recall:} say \emph{vīgintī annōs nātus sum} — \textbf{R4}, recognition at distance, eighty lessons back or more
  \item \emph{Recall:} say \emph{tempestās} — \textbf{R4}, recognition at distance, a second word from that far back
\end{itemize}
\end{quote}

\begin{cousinweb}
\textbf{quandō} — \textquotedblleft{}\textbf{when}.\textquotedblright{}

Built on the same \emph{\textbf{*kwo-}} stem as \emph{quī} and \emph{quid}, three lessons back. It does two jobs.

As a \textbf{question}: \textbf{Quandō venīs?} — \textquotedblleft{}When are you coming?\textquotedblright{}

As a \textbf{hinge}, joining two statements in time: \textbf{Quandō veniō, videō} — \textquotedblleft{}When I come, I see.\textquotedblright{}

Every Romance daughter kept it unchanged: Spanish \emph{cuándo}, Italian and Portuguese \emph{quando}, French \emph{quand}. English \textbf{quandary} is often attached to it and the connection is \textbf{not established} — the word appears in the sixteenth century from nowhere obvious, and \emph{quandō} is a guess rather than an etymology. This book will always tell you when a resemblance is only a resemblance.
\end{cousinweb}

\subsection*{Your turn}
Say these aloud:

\begin{itemize}
\raggedright
  \item \emph{quandō}
  \item every word this chapter has given you so far, in order
\end{itemize}

\begin{tcolorbox}[breakable,colback=teal!4,colframe=teal!35!black,title={Before you move on}]
Two jobs for \emph{quandō}? (\textbf{Asking when}, and \textbf{joining two statements in time}.) Is English \emph{quandary} from it? (\textbf{Not established} — a guess, not an etymology.)
\end{tcolorbox}
